% !TEX program = pdflatex
% Курс Deep RL от Hugging Face, модуль 3, русское учебное издание.
% Сборка из notes/ru/: latexmk -pdf unit3ru.tex; uses shared preamble.tex.
\PassOptionsToPackage{T2A}{fontenc}
\PassOptionsToPackage{russian}{babel}
\documentclass[11pt,a4paper]{article}

% Shared typography/layout for Hugging Face Deep RL Course notes.
% Compile with pdfLaTeX / latexmk -pdf with notes/ as the working directory.
% e.g. cd notes && latexmk -pdf unit1.tex

\usepackage[T1]{fontenc}
\usepackage[utf8]{inputenc}
\usepackage[english]{babel}
\usepackage{lmodern}
\usepackage{textcomp}
\usepackage{geometry}
\geometry{a4paper,margin=22mm,headheight=14pt,headsep=8mm,footskip=12mm}

\usepackage{amsmath,amssymb,mathtools}
\usepackage{graphicx,caption,needspace,placeins}
\usepackage{array,booktabs,tabularx}
\usepackage{xcolor}
\usepackage{microtype}
\usepackage{enumitem}
\usepackage{listings}
\usepackage[most]{tcolorbox}
\usepackage{titlesec}
\usepackage{fancyhdr}
\usepackage{hyperref}
\usepackage{bookmark}

\definecolor{ink}{HTML}{243041}
\definecolor{accent}{HTML}{38598C}
\definecolor{secondary}{HTML}{60768A}
\definecolor{codebg}{HTML}{F4F6F8}
\definecolor{codeborder}{HTML}{E3E8EE}
\definecolor{notebg}{HTML}{F3F7FC}

\hypersetup{
  colorlinks=true,
  linkcolor=accent,
  urlcolor=accent,
  citecolor=accent,
  pdfborder={0 0 0}
}
\urlstyle{same}

\setlength{\parindent}{0pt}
\setlength{\parskip}{0.38em plus 0.12em minus 0.05em}

\setlist[itemize]{
  topsep=0.35em,
  itemsep=0.13em,
  leftmargin=1.7em
}
\setlist[enumerate]{
  topsep=0.35em,
  itemsep=0.13em,
  leftmargin=1.8em
}

\setcounter{tocdepth}{2}
\setcounter{secnumdepth}{2}

\titleformat{\section}
  {\Large\bfseries\color{ink}}
  {\thesection}{0.7em}{}

\titleformat{\subsection}
  {\large\bfseries\color{accent}}
  {\thesubsection}{0.7em}{}

\titleformat{\subsubsection}
  {\normalsize\bfseries\color{ink}}
  {\thesubsubsection}{0.6em}{}

\titlespacing*{\section}
  {0pt}{2.2ex plus 0.7ex}{0.9ex}
\titlespacing*{\subsection}
  {0pt}{1.7ex plus 0.4ex}{0.6ex}
\titlespacing*{\subsubsection}
  {0pt}{1.3ex plus 0.3ex}{0.45ex}

\pagestyle{fancy}
\fancyhf{}

\fancyhead[L]{
  \footnotesize\sffamily\color{secondary}
  HF Deep RL Course \enspace/\enspace Unit 1
}
\fancyhead[R]{
  \footnotesize\sffamily\color{secondary}
  \nouppercase{\rightmark}
}
\fancyfoot[C]{
  \footnotesize\color{secondary}\thepage
}
\renewcommand{\headrulewidth}{0.25pt}

\captionsetup{
  font=small,
  labelfont={bf,color=secondary},
  textfont={color=ink},
  justification=centering,
  skip=4pt,
  hypcap=false
}

\lstdefinestyle{hfcode}{
  basicstyle=\ttfamily\footnotesize,
  backgroundcolor=\color{codebg},
  frame=single,
  rulecolor=\color{codeborder},
  numbers=none,
  showstringspaces=false,
  breaklines=true,
  breakatwhitespace=false,
  columns=fullflexible,
  keepspaces=true,
  xleftmargin=0pt,
  xrightmargin=0pt,
  aboveskip=0.6em,
  belowskip=0.6em,
  framexleftmargin=7pt,
  framexrightmargin=7pt,
  framextopmargin=5pt,
  framexbottommargin=5pt,
  upquote=true,
  tabsize=4,
}
\lstset{style=hfcode}

\newcommand{\CourseFigure}[2]{%
  \par\addvspace{0.55em}\noindent%
  \begin{minipage}{\linewidth}%
    \centering
    \includegraphics[
      width=0.86\linewidth,
      height=0.295\textheight,
      keepaspectratio
    ]{#1}\par
    \captionof{figure}{#2}%
  \end{minipage}%
  \par\addvspace{0.55em}%
}

\newtcolorbox{coursenote}{
  enhanced,
  breakable,
  colback=notebg,
  colframe=accent!30,
  boxrule=0.35pt,
  arc=1.2mm,
  left=9pt,
  right=9pt,
  top=6pt,
  bottom=6pt
}

\renewenvironment{quote}
  {\begin{coursenote}}
  {\end{coursenote}}

\emergencystretch=2.5em

\providecommand{\tightlist}{
  \setlength{\itemsep}{0pt}
  \setlength{\parskip}{0pt}
}

\providecommand{\passthrough}[1]{#1}

\Urlmuskip=0mu plus 6mu\relax

\renewcommand{\rmdefault}{cmr}
\renewcommand{\sfdefault}{cmss}
\renewcommand{\ttdefault}{cmtt}
\graphicspath{{../assets/unit3/}}
% The shared preamble has a Unit 1 label: override the unit-specific running header.
\fancyhead[L]{\footnotesize\sffamily\color{secondary}Курс Deep RL от Hugging Face \enspace/\enspace Модуль 3}
\fancyhead[R]{\footnotesize\sffamily\color{secondary}Глубокое Q-обучение}
\title{\textbf{Глубокое обучение с подкреплением}\\[0.35em]\large Модуль 3: глубокое Q-обучение}
\author{Курс Deep RL от Hugging Face\\\small Полное учебное издание на основе исходных MDX-файлов}
\date{}
\begin{document}
\selectlanguage{russian}
\fontencoding{T2A}\selectfont
\maketitle
\vspace{-0.8em}
\begin{center}\small\color{secondary}Исходный материал: курс Deep RL от Hugging Face.\end{center}
\tableofcontents
\vspace{1em}
\clearpage

\section{Глубокое Q-обучение}\label{deep-q-learning}

\CourseFigure{thumbnail.jpg}{Обложка модуля 3: глубокое Q-обучение \protect\href{https://huggingface.co/datasets/huggingface-deep-rl-course/course-images/resolve/main/en/unit4/thumbnail.jpg}{[оригинал]}}

В предыдущем модуле мы познакомились с первым алгоритмом обучения с подкреплением --- Q-обучением, \textbf{реализовали его с нуля} и обучили агента в двух средах: FrozenLake-v1 (замёрзшее озеро) и Taxi-v3 (такси).

С помощью этого простого алгоритма мы получили отличные результаты. Однако обе среды были относительно простыми, поскольку \textbf{пространство состояний в них было небольшим и дискретным}: 16 состояний в FrozenLake-v1 и 500 в Taxi-v3. Для сравнения, пространство состояний в играх Atari может \textbf{содержать от \(10^{9}\) до \(10^{11}\) состояний}.

Но, как мы увидим дальше, \textbf{создавать и обновлять Q-таблицу в средах с большим пространством состояний становится неэффективно}.

Поэтому в этом модуле мы \textbf{изучим первого агента глубокого обучения с подкреплением} --- Deep Q-Learning. Вместо Q-таблицы он использует нейронную сеть, которая получает состояние на вход и аппроксимирует Q-значения всех возможных действий для этого состояния.

Кроме того, \textbf{мы обучим агента играть в Space Invaders и другие игры Atari с помощью \href{https://github.com/DLR-RM/rl-baselines3-zoo}{RL-Zoo}}. Это инструментарий обучения с подкреплением на базе Stable-Baselines, содержащий скрипты для обучения и оценки агентов, подбора гиперпараметров, построения графиков и записи видео.

\CourseFigure{atari-envs.png}{Примеры сред Atari, в которых можно обучать агентов \protect\href{https://huggingface.co/datasets/huggingface-deep-rl-course/course-images/resolve/main/en/unit4/atari-envs.gif}{[оригинал]}}

Приступим!

\section{От Q-обучения к глубокому Q-обучению}\label{from-q-learning-to-deep-q-learning}

Мы узнали, что \textbf{Q-обучение --- это алгоритм для обучения Q-функции}, то есть \textbf{функции ценности действия (action-value function)}, которая определяет ценность выбора конкретного действия в конкретном состоянии.

\CourseFigure{Q-function.jpg}{Табличная Q-функция: значения для пар «состояние--действие» \protect\href{https://huggingface.co/datasets/huggingface-deep-rl-course/course-images/resolve/main/en/unit3/Q-function.jpg}{[оригинал]}}

Название \textbf{Q происходит от слова «качество» (quality) действия в данном состоянии}.

Внутри нашего алгоритма Q-функция представлена \textbf{Q-таблицей, каждая ячейка которой соответствует ценности некоторой пары «состояние--действие»}. Можно представить эту таблицу как \textbf{память или шпаргалку Q-функции}.

Проблема заключается в том, что Q-обучение --- \emph{табличный метод}. Если пространства состояний и действий \textbf{слишком велики, чтобы эффективно представить их массивами и таблицами}, такой метод перестаёт быть практичным. Иными словами, он \textbf{плохо масштабируется}. Q-обучение хорошо работало в небольших средах, например:

\begin{itemize}
\tightlist
\item
В FrozenLake было 16 состояний.
\item
В Taxi-v3 было 500 состояний.
\end{itemize}

Но сегодня мы хотим научить агента играть в Space Invaders --- более сложную игру, где входными данными служат изображения отдельных кадров.

Как \href{https://twitter.com/meln1k}{отмечал Никита Мелькозеров}, \textbf{пространство наблюдений Atari} имеет размерность (210, 160, 3)*. Значения компонентов находятся в диапазоне от 0 до 255; таким образом, всего возможно \(256^{210 \times 160 \times 3} = 256^{100800}\) разных наблюдений. Для сравнения: число атомов в наблюдаемой Вселенной оценивается примерно в \(10^{80}\).

\begin{itemize}
\tightlist
\item
Один кадр Atari --- это изображение размером 210 на 160 пикселей. Поскольку оно цветное (RGB), используются три канала; отсюда размерность (210, 160, 3). Значение каждого канала каждого пикселя может изменяться от 0 до 255.
\end{itemize}

\CourseFigure{atari.jpg}{Огромное пространство наблюдений, представленное изображениями Atari \protect\href{https://huggingface.co/datasets/huggingface-deep-rl-course/course-images/resolve/main/en/unit4/atari.jpg}{[оригинал]}}

Следовательно, пространство состояний чрезвычайно велико, и создавать и обновлять Q-таблицу для такой среды неэффективно. Вместо этого разумнее аппроксимировать Q-значения параметризованной Q-функцией \(Q_{\theta}(s,a)\).

Получая состояние на вход, нейронная сеть будет оценивать Q-значения всех допустимых действий в этом состоянии. Именно так и работает глубокое Q-обучение.

\CourseFigure{deep.jpg}{Нейронная сеть аппроксимирует ценности действий вместо использования Q-таблицы \protect\href{https://huggingface.co/datasets/huggingface-deep-rl-course/course-images/resolve/main/en/unit1/deep.jpg}{[оригинал]}}

Теперь, разобравшись в общей идее Deep Q-Learning, рассмотрим подробнее архитектуру Deep Q-Network.

\section{Глубокая Q-сеть (DQN)}\label{the-deep-q-network-dqn}

Вот как устроена нейронная сеть для глубокого Q-обучения:

\CourseFigure{deep-q-network.jpg}{Архитектура DQN: стек последовательных кадров и оценки ценности действий \protect\href{https://huggingface.co/datasets/huggingface-deep-rl-course/course-images/resolve/main/en/unit4/deep-q-network.jpg}{[оригинал]}}

На вход сети подаётся \textbf{стек из четырёх кадров}, представляющий состояние, а на выходе получается \textbf{вектор Q-значений для всех возможных действий в этом состоянии}. Затем, как и в обычном Q-обучении, мы выбираем действие с помощью \(\varepsilon\)-жадной стратегии.

При инициализации нейронной сети \textbf{оценки Q-значений крайне неточны}. Однако в ходе обучения агент DQN начинает связывать игровые ситуации с подходящими действиями и \textbf{учится успешно играть}.

\subsection{Предобработка входа и проблема временной информации}\label{preprocessing-the-input-and-temporal-limitation}

Необходимо \textbf{предобрабатывать входные данные}. Это важный этап: мы стремимся \textbf{уменьшить сложность представления состояния и сократить вычислительные затраты на обучение}.

Для этого мы \textbf{уменьшаем размер изображения до 84 на 84 пикселя и переводим его в оттенки серого}. Цвета в играх Atari обычно не несут критически важной информации, поэтому мы можем \textbf{заменить три канала RGB одним каналом} и существенно сократить объём входных данных.

В некоторых играх можно также \textbf{обрезать неинформативную часть экрана}. После этого мы объединяем четыре последовательных кадра в один стек.

\CourseFigure{preprocessing.jpg}{Предобработка наблюдений Atari и формирование стека кадров \protect\href{https://huggingface.co/datasets/huggingface-deep-rl-course/course-images/resolve/main/en/unit4/preprocessing.jpg}{[оригинал]}}

\textbf{Зачем объединять четыре кадра?} Это позволяет \textbf{устранить недостаток временной информации}. Рассмотрим пример с игрой Pong. Допустим, мы видим такой кадр:

\CourseFigure{temporal-limitation.jpg}{По одному кадру Pong нельзя определить скорость и направление движения мяча \protect\href{https://huggingface.co/datasets/huggingface-deep-rl-course/course-images/resolve/main/en/unit4/temporal-limitation.jpg}{[оригинал]}}

Можете ли вы сказать, куда движется мяч? Нет: одного кадра недостаточно, чтобы понять движение! А если добавить ещё три кадра? \textbf{Теперь видно, что мяч движется вправо}.

\CourseFigure{temporal-limitation-2.jpg}{Несколько последовательных кадров Pong позволяют определить направление мяча \protect\href{https://huggingface.co/datasets/huggingface-deep-rl-course/course-images/resolve/main/en/unit4/temporal-limitation-2.jpg}{[оригинал]}}

Именно поэтому для получения информации о движении мы объединяем четыре кадра.

Далее стек кадров проходит через три свёрточных слоя. Эти слои \textbf{извлекают и используют пространственные зависимости на изображениях}. Поскольку кадры объединены, сеть также \textbf{может использовать некоторые временные зависимости между ними}.

Если вы пока не знаете, что такое свёрточные слои, не переживайте. Можно изучить \href{https://www.udacity.com/course/deep-learning-pytorch--ud188}{четвёртый урок бесплатного курса по глубокому обучению от Udacity}.

Наконец, несколько полносвязных слоёв формируют Q-значения для всех возможных действий в данном состоянии.

\CourseFigure{deep-q-network.jpg}{Архитектура DQN: стек последовательных кадров и оценки ценности действий \protect\href{https://huggingface.co/datasets/huggingface-deep-rl-course/course-images/resolve/main/en/unit4/deep-q-network.jpg}{[оригинал]}}

Итак, Deep Q-Learning использует нейронную сеть, которая по состоянию аппроксимирует Q-значения всех доступных действий. Теперь изучим сам алгоритм обучения.

\section{Алгоритм глубокого Q-обучения}\label{the-deep-q-learning-algorithm}

Мы выяснили, что глубокое Q-обучение \textbf{использует глубокую нейронную сеть для аппроксимации Q-значений всех возможных действий в каждом состоянии} (оценки функции ценности).

Отличие заключается в том, что во время обучения мы не обновляем Q-значение пары «состояние--действие» непосредственно, как в обычном Q-обучении:

\CourseFigure{q-ex-5.jpg}{Обновление в табличном Q-обучении как отправная точка \protect\href{https://huggingface.co/datasets/huggingface-deep-rl-course/course-images/resolve/main/en/unit3/q-ex-5.jpg}{[оригинал]}}

в Deep Q-Learning мы определяем \textbf{функцию потерь, сравнивающую предсказанное Q-значение с целевым Q-значением. Затем с помощью градиентного спуска обновляем веса глубокой Q-сети}, чтобы она лучше аппроксимировала Q-значения.

\CourseFigure{Q-target.jpg}{Функция потерь Deep Q-Learning и целевое значение TD \protect\href{https://huggingface.co/datasets/huggingface-deep-rl-course/course-images/resolve/main/en/unit4/Q-target.jpg}{[оригинал]}}

Алгоритм обучения Deep Q-Learning включает \emph{две фазы}:

\begin{itemize}
\tightlist
\item
\textbf{Сбор опыта (sampling):} агент совершает действия и \textbf{сохраняет наблюдаемые переходы в буфере воспроизведения опыта (replay memory)}.
\item
\textbf{Обучение (training):} мы случайным образом выбираем \textbf{небольшой пакет переходов и обучаемся на нём, выполняя шаг обновления весов градиентным спуском}.
\end{itemize}

\CourseFigure{sampling-training.jpg}{Две фазы Deep Q-Learning: сбор опыта и обучение \protect\href{https://huggingface.co/datasets/huggingface-deep-rl-course/course-images/resolve/main/en/unit4/sampling-training.jpg}{[оригинал]}}

Но это не единственное отличие от табличного Q-обучения. Обучение Deep Q-Learning \textbf{может быть нестабильным}, главным образом из-за сочетания нелинейной функции Q-значений (нейронной сети) и бутстрэппинга (когда целевые значения обновляют на основе уже имеющихся оценок, а не фактически полученной полной суммы вознаграждений).

Чтобы стабилизировать обучение, применяют три разных приёма:

\begin{enumerate}
\def\labelenumi{\arabic{enumi}.}
\tightlist
\item
\emph{Воспроизведение опыта (Experience Replay)} --- для \textbf{более эффективного использования накопленных переходов}.
\item
\emph{Фиксированная целевая Q-сеть (Fixed Q-Target)} --- \textbf{для стабилизации обучения}.
\item
\emph{Двойное глубокое Q-обучение (Double DQN)} --- \textbf{для уменьшения переоценки Q-значений}.
\end{enumerate}

Рассмотрим каждый из них!

\subsection{Воспроизведение опыта для более эффективного обучения}\label{experience-replay-to-make-more-efficient-use-of-experiences}

Зачем нам нужна память воспроизведения опыта?

В Deep Q-Learning механизм Experience Replay выполняет две функции:

\begin{enumerate}
\def\labelenumi{\arabic{enumi}.}
\tightlist
\item
\textbf{Более эффективное использование опыта во время обучения}. При обычном онлайн-обучении с подкреплением агент взаимодействует со средой, получает переходы (состояние, действие, вознаграждение и следующее состояние), учится на них (обновляет нейронную сеть), а затем отбрасывает. Это неэффективно.
\end{enumerate}

Experience Replay позволяет \textbf{эффективнее использовать уже полученные данные}. Мы храним переходы в буфере воспроизведения и \textbf{повторно используем их на протяжении обучения}.

\CourseFigure{experience-replay.jpg}{Буфер воспроизведения опыта, сохраняющий предыдущие переходы \protect\href{https://huggingface.co/datasets/huggingface-deep-rl-course/course-images/resolve/main/en/unit4/experience-replay.jpg}{[оригинал]}}

=\textgreater{} Благодаря этому агент может \textbf{обучаться на одних и тех же переходах многократно}.

\begin{enumerate}
\def\labelenumi{\arabic{enumi}.}
\setcounter{enumi}{1}
\tightlist
\item
\textbf{Предотвращение забывания предыдущего опыта (катастрофической интерференции) и уменьшение корреляции между последовательными переходами}.
\end{enumerate}

\begin{itemize}
\tightlist
\item
\textbf{\href{https://en.wikipedia.org/wiki/Catastrophic_interference}{Катастрофическое забывание}}: если подавать нейронной сети последовательные переходы, она склонна \textbf{забывать предыдущий опыт по мере поступления нового}. Например, когда агент переходит с первого уровня игры на сильно отличающийся второй, он может забыть, как успешно проходить первый уровень.
\end{itemize}

Решение --- сохранять переходы при взаимодействии со средой в буфере Replay Buffer, а затем выбирать оттуда небольшие случайные пакеты. Это не позволяет \textbf{нейросети учиться исключительно на том, что агент делал непосредственно перед обновлением}.

У воспроизведения опыта есть и другие преимущества. Случайная выборка уменьшает корреляции между последовательными наблюдениями и помогает предотвратить \textbf{сильные колебания или расходимость оценок ценности действий}.

В псевдокоде Deep Q-Learning мы \textbf{создаём буфер памяти \(D\) вместимостью \(N\)} (\(N\) --- задаваемый гиперпараметр). Затем сохраняем в нём переходы и во время обучения выбираем из буфера пакет опыта для передачи в глубокую Q-сеть.

\CourseFigure{experience-replay-pseudocode.jpg}{Псевдокод DQN: использование буфера воспроизведения опыта \protect\href{https://huggingface.co/datasets/huggingface-deep-rl-course/course-images/resolve/main/en/unit4/experience-replay-pseudocode.jpg}{[оригинал]}}

\subsection{Фиксированная целевая Q-сеть для стабилизации обучения}\label{fixed-q-target-to-stabilize-the-training}

При вычислении ошибки временных разностей (TD error, то есть функции потерь) мы находим \textbf{разность между целевым TD-значением (Q-Target) и текущей оценкой Q-значения}.

Однако \textbf{истинное целевое TD-значение нам неизвестно}, и его приходится оценивать. Из уравнения Беллмана следует, что цель равна вознаграждению за действие в текущем состоянии плюс дисконтированное максимальное Q-значение в следующем состоянии.

\CourseFigure{Q-target.jpg}{Функция потерь Deep Q-Learning и целевое значение TD \protect\href{https://huggingface.co/datasets/huggingface-deep-rl-course/course-images/resolve/main/en/unit4/Q-target.jpg}{[оригинал]}}

Проблема в том, что для вычисления целевого TD-значения \textbf{и} текущей оценки Q мы используем одни и те же параметры (веса). Поэтому между целевым значением и изменяемыми параметрами возникает сильная взаимосвязь.

Следовательно, на каждом шаге обучения \textbf{смещаются и наши Q-оценки, и целевые значения}. Мы приближаемся к цели, но она сама постоянно движется. Это похоже на погоню за движущейся мишенью и может вызывать значительные колебания процесса обучения.

Представьте, что вы ковбой (текущая Q-оценка) и хотите поймать корову (целевое Q-значение). Ваша задача --- приблизиться к ней (уменьшить ошибку).

\CourseFigure{qtarget-1.jpg}{Аналогия с движущейся целью: первый шаг обучения Q-функции \protect\href{https://huggingface.co/datasets/huggingface-deep-rl-course/course-images/resolve/main/en/unit4/qtarget-1.jpg}{[оригинал]}}

На каждом временном шаге вы пытаетесь приблизиться к корове, но она тоже перемещается на каждом шаге (потому что используются одни и те же параметры).

\CourseFigure{qtarget-2.jpg}{Движущаяся цель: смещение целевого значения \protect\href{https://huggingface.co/datasets/huggingface-deep-rl-course/course-images/resolve/main/en/unit4/qtarget-2.jpg}{[оригинал]}}

\CourseFigure{qtarget-3.jpg}{Движущаяся цель: последовательные изменения положения цели \protect\href{https://huggingface.co/datasets/huggingface-deep-rl-course/course-images/resolve/main/en/unit4/qtarget-3.jpg}{[оригинал]}}

В результате траектория преследования становится странной и извилистой, что соответствует сильным колебаниям во время обучения.

\CourseFigure{qtarget-4.jpg}{Движущаяся цель: колебания при обучении функции ценности \protect\href{https://huggingface.co/datasets/huggingface-deep-rl-course/course-images/resolve/main/en/unit4/qtarget-4.jpg}{[оригинал]}}

Вместо этого, как видно из псевдокода, мы:

\begin{itemize}
\tightlist
\item
Используем \textbf{отдельную сеть с фиксированными параметрами} для оценки целевого TD-значения.
\item
\textbf{Копируем параметры основной глубокой Q-сети каждые \(C\) шагов}, обновляя целевую сеть.
\end{itemize}

\CourseFigure{fixed-q-target-pseudocode.jpg}{Псевдокод DQN: обновление фиксированной целевой сети \protect\href{https://huggingface.co/datasets/huggingface-deep-rl-course/course-images/resolve/main/en/unit4/fixed-q-target-pseudocode.jpg}{[оригинал]}}

\subsection{Двойная DQN (Double DQN)}\label{double-dqn}

Алгоритм Double DQN, или двойное глубокое Q-обучение, предложил \href{https://papers.nips.cc/paper/3964-double-q-learning}{Хадо ван Хасселт}. Этот метод \textbf{решает проблему переоценки Q-значений}.

Чтобы понять проблему, вспомним, как вычисляется целевое TD-значение:

\CourseFigure{TD-1.jpg}{Цель TD и операция максимизации по Q-значениям \protect\href{https://huggingface.co/datasets/huggingface-deep-rl-course/course-images/resolve/main/en/unit3/TD-1.jpg}{[оригинал]}}

При вычислении целевого TD-значения возникает простой вопрос: откуда мы знаем, что \textbf{лучшее действие для следующего состояния --- это именно действие с максимальным Q-значением}?

Нам известно, что точность Q-оценок зависит и от того, \textbf{какие действия мы уже пробовали}, и от того, какие соседние состояния исследовали.

Поэтому в начале обучения информации о действительно лучшем действии недостаточно. Если считать максимум зашумлённых Q-оценок лучшим действием, можно получить ложноположительные результаты. Когда неоптимальным действиям регулярно \textbf{приписываются более высокие Q-значения, чем действительно оптимальному действию}, обучение осложняется.

Решение: при вычислении целевого Q-значения использовать две сети, чтобы разделить выбор действия и получение целевой оценки ценности. Для этого мы:

\begin{itemize}
\tightlist
\item
Используем \textbf{основную сеть DQN}, чтобы выбрать лучшее действие для следующего состояния (действие с наибольшим Q-значением).
\item
Используем \textbf{целевую сеть}, чтобы рассчитать целевое Q-значение выбранного действия в следующем состоянии.
\end{itemize}

Таким образом, Double DQN помогает уменьшить переоценку Q-значений и, как следствие, обеспечивает более быстрое и устойчивое обучение.

Помимо трёх рассмотренных усовершенствований Deep Q-Learning существуют и другие, например Prioritized Experience Replay и Dueling Deep Q-Learning. Они выходят за рамки этого курса, но при желании вы можете изучить материалы по ссылкам из списка дополнительной литературы.

\section{Практическое занятие}\label{hands-on}

\href{https://colab.research.google.com/github/huggingface/deep-rl-class/blob/main/notebooks/unit3/unit3.ipynb}{Открыть официальный ноутбук модуля 3 в Google Colab}

\href{http://hf.co/join/discord}{Официальная помощь по курсу и Discord}

Теперь, когда мы изучили теорию Deep Q-Learning, \textbf{вы готовы обучить агента играть в игры Atari}. Начнём со Space Invaders, но затем сможете выбрать любую другую игру Atari.

\CourseFigure{atari-envs.png}{Примеры сред Atari, в которых можно обучать агентов \protect\href{https://huggingface.co/datasets/huggingface-deep-rl-course/course-images/resolve/main/en/unit4/atari-envs.gif}{[оригинал]}}

Мы используем \href{https://github.com/DLR-RM/rl-baselines3-zoo}{интеграцию RL-Baselines3 Zoo}: обычную реализацию Deep Q-Learning без расширений Double-DQN, Dueling-DQN и Prioritized Experience Replay.

Если после практики вы захотите \textbf{реализовать Deep Q-Learning самостоятельно}, обязательно посмотрите реализацию CleanRL: \url{https://github.com/vwxyzjn/cleanrl/blob/master/cleanrl/dqn_atari.py}

Чтобы практическое задание засчиталось для сертификации, необходимо загрузить обученную модель в Hub и \textbf{получить результат не ниже 200}.

Чтобы узнать результат, найдите модель в таблице лидеров: \textbf{результат = mean\_reward - std вознаграждения}.

\textbf{Если модели нет в списке, прокрутите страницу вниз и нажмите кнопку обновления.}

Подробнее о сертификации см. в \href{https://huggingface.co/deep-rl-course/en/unit0/introduction#certification-process}{соответствующем разделе курса}.

А здесь можно \href{https://huggingface.co/spaces/ThomasSimonini/Check-my-progress-Deep-RL-Course}{проверить собственный прогресс}.

\textbf{Чтобы начать практику, нажмите кнопку «Open In Colab»:}

\href{https://colab.research.google.com/github/huggingface/deep-rl-class/blob/master/notebooks/unit3/unit3.ipynb}{Открыть в Colab}

\subsection{Модуль 3: Deep Q-Learning в играх Atari (Space Invaders) с помощью RL Baselines3 Zoo}\label{unit-3-deep-q-learning-with-atari-games-space-invaders-using-rl-baselines3-zoo}

\CourseFigure{thumbnail.jpg}{Обложка модуля 3: глубокое Q-обучение \protect\href{https://huggingface.co/datasets/huggingface-deep-rl-course/course-images/resolve/main/en/unit4/thumbnail.jpg}{[оригинал]}}

В этой практической работе \textbf{вы обучите агента Deep Q-Learning играть в Space Invaders} с помощью \href{https://github.com/DLR-RM/rl-baselines3-zoo}{RL Baselines3 Zoo} --- инструментария на основе \href{https://stable-baselines3.readthedocs.io/en/master/}{Stable-Baselines3}, содержащего скрипты для обучения и оценки агентов, настройки гиперпараметров, построения графиков и записи видео.

Мы используем \href{https://stable-baselines3.readthedocs.io/en/master/modules/dqn.html}{интеграцию RL-Baselines3 Zoo с базовой реализацией Deep Q-Learning} без таких расширений, как Double-DQN, Dueling-DQN и Prioritized Experience Replay.

\paragraph{Игровые среды}\label{game-environments}

\begin{itemize}
\tightlist
\item
\href{https://gymnasium.farama.org/environments/atari/space_invaders/}{SpacesInvadersNoFrameskip-v4}
\end{itemize}

Различия между версиями Space Invaders описаны \href{https://gymnasium.farama.org/environments/atari/space_invaders/#variants}{здесь}.

\paragraph{Библиотека RL для изучения}\label{reading-rl-library}

\begin{itemize}
\tightlist
\item
\href{https://github.com/DLR-RM/rl-baselines3-zoo}{RL-Baselines3-Zoo}
\end{itemize}

\subsubsection{Цели практического занятия}\label{objectives-of-this-hands-on}

После выполнения работы вы сможете:

\begin{itemize}
\tightlist
\item
Глубже понять, \textbf{как работает RL Baselines3 Zoo}.
\item
\textbf{Загрузить обученного агента и код в Hugging Face Hub} вместе с красивой видеозаписью игры и результатами оценки.
\end{itemize}

\subsubsection{Предварительные требования}\label{prerequisites}

Перед выполнением практического задания необходимо:

[Чтение] \textbf{\href{https://huggingface.co/deep-rl-course/unit3/introduction}{Изучить теорию Deep Q-Learning из модуля 3}}.

Мы постоянно улучшаем наши руководства. Поэтому, \textbf{если вы заметили проблему в практическом задании}, пожалуйста, \href{https://github.com/huggingface/deep-rl-class/issues}{создайте issue в репозитории GitHub}.

\subsection{Обучим агента Deep Q-Learning играть в Atari Space Invaders и загрузим его в Hub}\label{lets-train-a-deep-q-learning-agent-playing-atari-space-invaders-space-invaders-and-upload-it-to-the-hub.}

Мы настоятельно рекомендуем студентам \textbf{выполнять практические задания в Google Colab, а не на личных компьютерах}.

Используя Google Colab, \textbf{можно сосредоточиться на обучении и экспериментах, не отвлекаясь на технические вопросы настройки окружения}.

Чтобы практическое задание засчиталось для сертификации, нужно загрузить обученную модель в Hub и \textbf{получить результат не ниже 200}.

Чтобы найти результат, откройте таблицу лидеров и найдите модель. \textbf{Результат вычисляется как mean\_reward минус стандартное отклонение вознаграждения}.

Подробнее о сертификации см. \href{https://huggingface.co/deep-rl-course/en/unit0/introduction#certification-process}{в разделе курса}.

\subsubsection{Настройка GPU}\label{set-the-gpu}

\begin{itemize}
\tightlist
\item
Чтобы \textbf{ускорить обучение агента, воспользуемся GPU}. В Colab выберите \texttt{Runtime\ \textgreater{}\ Change\ Runtime\ type}.
\end{itemize}

\CourseFigure{gpu-step1.jpg}{Google Colab: изменение типа среды выполнения \protect\href{https://huggingface.co/datasets/huggingface-deep-rl-course/course-images/resolve/main/en/notebooks/gpu-step1.jpg}{[оригинал]}}

\begin{itemize}
\tightlist
\item
\texttt{Hardware\ Accelerator\ \textgreater{}\ GPU}
\end{itemize}

\CourseFigure{gpu-step2.jpg}{Google Colab: выбор ускорителя GPU \protect\href{https://huggingface.co/datasets/huggingface-deep-rl-course/course-images/resolve/main/en/notebooks/gpu-step2.jpg}{[оригинал]}}

\subsection{Установка RL-Baselines3 Zoo и зависимостей (материал для чтения)}\label{install-rl-baselines3-zoo-and-its-dependencies-reading}

Если появится сообщение \texttt{ERROR:\ pip\textquotesingle{}s\ dependency\ resolver\ does\ not\ currently\ take\ into\ account\ all\ the\ packages\ that\ are\ installed.}, \textbf{это нормально и не является критической ошибкой}: возник конфликт версий, но необходимые нам пакеты установлены.

\begin{lstlisting}
# For now we install this update of RL-Baselines3 Zoo
pip install git+https://github.com/DLR-RM/rl-baselines3-zoo
\end{lstlisting}

\begin{lstlisting}
apt-get install swig cmake ffmpeg
\end{lstlisting}

Для игр Atari в Gymnasium необходимо установить пакет поддержки Atari. Кроме того, нужно принять лицензию ROM-файлов с играми с помощью пакета \texttt{accept-rom-license}.

\begin{lstlisting}
!pip install gymnasium[atari]
!pip install gymnasium[accept-rom-license]
\end{lstlisting}

\subsubsection{Создание виртуального дисплея}\label{create-a-virtual-display}

Во время практики понадобится записать видео игры. Если обучение выполняется на компьютере без монитора (headless), \textbf{необходимо создать виртуальный экран, чтобы среда могла отрисовывать кадры} и их можно было записать.

Следующая ячейка устанавливает библиотеки, создаёт виртуальный экран и запускает его.

\begin{lstlisting}
apt install python-opengl
apt install ffmpeg
apt install xvfb
pip3 install pyvirtualdisplay
\end{lstlisting}

\begin{lstlisting}
# Virtual display
from pyvirtualdisplay import Display

virtual_display = Display(visible=0, size=(1400, 900))
virtual_display.start()
\end{lstlisting}

\subsubsection{Обучение агента Deep Q-Learning игре Space Invaders}\label{train-our-deep-q-learning-agent-to-play-space-invaders-space-invaders}

Чтобы обучить агента с помощью RL-Baselines3-Zoo, достаточно сделать две вещи:

\begin{enumerate}
\def\labelenumi{\arabic{enumi}.}
\tightlist
\item
Создать файл конфигурации гиперпараметров обучения с именем \texttt{dqn.yml}.
\end{enumerate}

Вот пример шаблона:

\begin{lstlisting}
SpaceInvadersNoFrameskip-v4:
  env_wrapper:
    - stable_baselines3.common.atari_wrappers.AtariWrapper
  frame_stack: 4
  policy: 'CnnPolicy'
  n_timesteps: !!float 1e7
  buffer_size: 100000
  learning_rate: !!float 1e-4
  batch_size: 32
  learning_starts: 100000
  target_update_interval: 1000
  train_freq: 4
  gradient_steps: 1
  exploration_fraction: 0.1
  exploration_final_eps: 0.01
  # If True, you need to deactivate handle_timeout_termination
  # in the replay_buffer_kwargs
  optimize_memory_usage: False
\end{lstlisting}

Здесь можно увидеть следующее:

\begin{itemize}
\tightlist
\item
Мы используем \texttt{Atari\ Wrapper}, который предобрабатывает вход (уменьшает изображение, переводит в оттенки серого и объединяет четыре кадра).
\item
Мы используем \texttt{CnnPolicy}, поскольку обработку кадров выполняют свёрточные слои.
\item
Обучаем агента на протяжении 10 миллионов временных шагов (\texttt{n\_timesteps}).
\item
Размер буфера Experience Replay составляет 100000 --- столько переходов сохраняется для последующего повторного обучения агента.
\end{itemize}

Совет: рекомендую \textbf{уменьшить число шагов обучения до одного миллиона}. На GPU P100 это займёт около 90 минут. Команда \texttt{!nvidia-smi} покажет, какой GPU у вас используется. На 10 миллионов шагов потребуется около девяти часов. Такое обучение лучше выполнять на собственном компьютере или другой машине. Для скачивания ноутбука выберите \texttt{File\textgreater{}Download}.

При оптимизации гиперпараметров рекомендую сосредоточиться на трёх параметрах:

\begin{itemize}
\tightlist
\item
\texttt{learning\_rate}
\item
\texttt{buffer\_size} (размер памяти опыта)
\item
\texttt{batch\_size}
\end{itemize}

Полезная привычка --- \textbf{сверяться с документацией и понимать назначение каждого гиперпараметра}: \url{https://stable-baselines3.readthedocs.io/en/master/modules/dqn.html#parameters}

\begin{enumerate}
\def\labelenumi{\arabic{enumi}.}
\setcounter{enumi}{1}
\tightlist
\item
Начинаем обучение и сохраняем модели в папку \texttt{logs}.
\end{enumerate}

\begin{itemize}
\tightlist
\item
Указываем алгоритм после \texttt{-{}-algo}, папку сохранения модели после \texttt{-f}, а файл конфигурации гиперпараметров после \texttt{-c}.
\end{itemize}

\begin{lstlisting}
python -m rl_zoo3.train --algo ________ --env SpaceInvadersNoFrameskip-v4  -f _________  -c _________
\end{lstlisting}

\subparagraph{Решение}\label{solution}

\begin{lstlisting}
python -m rl_zoo3.train --algo dqn  --env SpaceInvadersNoFrameskip-v4 -f logs/ -c dqn.yml
\end{lstlisting}

\subsubsection{Оценим нашего агента}\label{lets-evaluate-our-agent}

\begin{itemize}
\tightlist
\item
RL-Baselines3-Zoo предоставляет \texttt{enjoy.py} --- Python-скрипт для оценки агента. Во многих библиотеках RL скрипт оценки называется именно \texttt{enjoy.py}.
\item
Оценим его на протяжении 5000 временных шагов.
\end{itemize}

\begin{lstlisting}
python -m rl_zoo3.enjoy  --algo dqn  --env SpaceInvadersNoFrameskip-v4  --no-render  --n-timesteps _________  --folder logs/
\end{lstlisting}

\subparagraph{Решение}\label{solution-1}

\begin{lstlisting}
python -m rl_zoo3.enjoy  --algo dqn  --env SpaceInvadersNoFrameskip-v4  --no-render  --n-timesteps 5000  --folder logs/
\end{lstlisting}

\subsubsection{Публикация обученной модели в Hub}\label{publish-our-trained-model-on-the-hub}

Убедившись, что после обучения получены хорошие результаты, мы можем опубликовать обученную модель в Hub одной строкой кода.

\CourseFigure{space-invaders-model.png}{Демонстрация обученного агента Space Invaders (первый кадр) \protect\href{https://huggingface.co/datasets/huggingface-deep-rl-course/course-images/resolve/main/en/notebooks/unit3/space-invaders-model.gif}{[оригинал]}}

Команда \texttt{rl\_zoo3.push\_to\_hub} \textbf{проводит оценку, записывает видео, создаёт карточку модели и загружает агента в Hub}.

Таким образом, вы сможете:

\begin{itemize}
\tightlist
\item
\textbf{Показать результаты своей работы}.
\item
\textbf{Посмотреть на игру обученного агента}.
\item
\textbf{Поделиться с сообществом агентом, которым смогут пользоваться другие}.
\item
\textbf{Сравнить результаты агента с результатами других участников в таблице лидеров}: \url{https://huggingface.co/spaces/huggingface-projects/Deep-Reinforcement-Learning-Leaderboard}
\end{itemize}

Для публикации модели в сообществе осталось выполнить ещё три шага:

\begin{enumerate}
\def\labelenumi{\arabic{enumi}.}
\item
Если аккаунта ещё нет, создайте его на Hugging Face: \url{https://huggingface.co/join}
\item
Войдите в аккаунт, затем сохраните токен аутентификации, полученный на сайте Hugging Face.
\end{enumerate}

\begin{itemize}
\tightlist
\item
Создайте новый токен (\url{https://huggingface.co/settings/tokens}) \textbf{с правом записи (write)}.
\end{itemize}

\CourseFigure{create-token.jpg}{Создание токена доступа Hugging Face \protect\href{https://huggingface.co/datasets/huggingface-deep-rl-course/course-images/resolve/main/en/notebooks/create-token.jpg}{[оригинал]}}

\begin{itemize}
\tightlist
\item
Скопируйте токен.
\item
Запустите следующую ячейку и вставьте токен.
\end{itemize}

\begin{lstlisting}
from huggingface_hub import notebook_login # To log to our Hugging Face account to be able to upload models to the Hub.
notebook_login()
!git config --global credential.helper store
\end{lstlisting}

Если вы не используете Google Colab или Jupyter Notebook, вместо этого необходимо выполнить команду \texttt{huggingface-cli\ login}.

\begin{enumerate}
\def\labelenumi{\arabic{enumi}.}
\setcounter{enumi}{2}
\tightlist
\item
Теперь всё готово для отправки обученного агента в Hub.
\end{enumerate}

Запустим файл \texttt{push\_to\_hub.py}, чтобы загрузить обученного агента в Hub.

\texttt{-{}-repo-name}: название репозитория.

\texttt{-orga}: ваше имя пользователя на Hugging Face.

\texttt{-f}: папка, в которой находится обученная модель (в нашем случае \texttt{logs}).

\CourseFigure{select-id.png}{Выбор каталога с обученной моделью для публикации \protect\href{https://huggingface.co/datasets/huggingface-deep-rl-course/course-images/resolve/main/en/notebooks/unit3/select-id.png}{[оригинал]}}

\begin{lstlisting}
python -m rl_zoo3.push_to_hub  --algo dqn  --env SpaceInvadersNoFrameskip-v4  --repo-name _____________________ -orga _____________________ -f logs/
\end{lstlisting}

\subparagraph{Решение}\label{solution-2}

\begin{lstlisting}
python -m rl_zoo3.push_to_hub  --algo dqn  --env SpaceInvadersNoFrameskip-v4  --repo-name dqn-SpaceInvadersNoFrameskip-v4  -orga ThomasSimonini  -f logs/
\end{lstlisting}

\begin{center}\rule{0.5\linewidth}{0.5pt}\end{center}

Поздравляем: вы только что обучили и загрузили своего первого агента Deep Q-Learning с помощью RL-Baselines3 Zoo! Скрипт выше должен был вывести ссылку на репозиторий модели, например \url{https://huggingface.co/ThomasSimonini/dqn-SpaceInvadersNoFrameskip-v4}. Перейдя по ссылке, вы сможете:

\begin{itemize}
\tightlist
\item
Увидеть \textbf{предпросмотр видео с игрой агента} справа.
\item
Нажать «Files and versions» и просмотреть все файлы репозитория.
\item
Нажать «Use in stable-baselines3», чтобы получить пример кода для загрузки модели.
\item
Просмотреть карточку модели (файл \texttt{README.md}) с её описанием и гиперпараметрами обучения.
\end{itemize}

В основе Hub лежат репозитории Git (не беспокойтесь, если вы ещё не знакомы с Git). Поэтому можно публиковать новые версии модели по мере проведения экспериментов и улучшения агента.

\textbf{Сравните результаты своих агентов с результатами других участников} в \href{https://huggingface.co/spaces/huggingface-projects/Deep-Reinforcement-Learning-Leaderboard}{таблице лидеров}.

\subsubsection{Загрузка хорошо обученной модели}\label{load-a-powerful-trained-model}

\begin{itemize}
\tightlist
\item
Команда Stable-Baselines3 опубликовала в Hub \textbf{более 150 обученных агентов глубокого обучения с подкреплением}.
\end{itemize}

Найти их можно здесь: \url{https://huggingface.co/sb3}

Несколько примеров:

\begin{itemize}
\tightlist
\item
Asteroids: \url{https://huggingface.co/sb3/dqn-AsteroidsNoFrameskip-v4}
\item
Beam Rider: \url{https://huggingface.co/sb3/dqn-BeamRiderNoFrameskip-v4}
\item
Breakout: \url{https://huggingface.co/sb3/dqn-BreakoutNoFrameskip-v4}
\item
Road Runner: \url{https://huggingface.co/sb3/dqn-RoadRunnerNoFrameskip-v4}
\end{itemize}

Загрузим агента, играющего в Beam Rider. Его модель размещена здесь: \url{https://huggingface.co/sb3/dqn-BeamRiderNoFrameskip-v4}

\begin{enumerate}
\def\labelenumi{\arabic{enumi}.}
\tightlist
\item
Скачаем модель с помощью \texttt{rl\_zoo3.load\_from\_hub} и поместим в новую папку, которую назовём \texttt{rl\_trained}.
\end{enumerate}

\begin{lstlisting}
# Download model and save it into the logs/ folder
python -m rl_zoo3.load_from_hub --algo dqn --env BeamRiderNoFrameskip-v4 -orga sb3 -f rl_trained/
\end{lstlisting}

\begin{enumerate}
\def\labelenumi{\arabic{enumi}.}
\setcounter{enumi}{1}
\tightlist
\item
Оценим модель на протяжении 5000 временных шагов.
\end{enumerate}

\begin{lstlisting}
python -m rl_zoo3.enjoy --algo dqn --env BeamRiderNoFrameskip-v4 -n 5000  -f rl_trained/ --no-render
\end{lstlisting}

Почему бы не попробовать самостоятельно обучить \textbf{агента Deep Q-Learning} для игры Beam Rider (\texttt{BeamRiderNoFrameskip-v4})?

Если захотите попробовать, изучите карточку обученной модели: \url{https://huggingface.co/sb3/dqn-BeamRiderNoFrameskip-v4#hyperparameters}\par\noindent \textbf{В ней приведены использованные гиперпараметры обучения}.

Однако подбор гиперпараметров бывает непростой задачей. К счастью, в следующем модуле мы увидим, как \textbf{использовать Optuna для оптимизации гиперпараметров}.

\subsubsection{Дополнительные задания}\label{some-additional-challenges}

Лучший способ учиться --- \textbf{пробовать всё самостоятельно}!

В \href{https://huggingface.co/spaces/huggingface-projects/Deep-Reinforcement-Learning-Leaderboard}{таблице лидеров} вы найдёте своих агентов. Сможете ли вы подняться на первую строчку?

Вот несколько сред, в которых можно попробовать обучить агента:

\begin{itemize}
\tightlist
\item
BeamRiderNoFrameskip-v4
\item
BreakoutNoFrameskip-v4
\item
EnduroNoFrameskip-v4
\item
PongNoFrameskip-v4
\end{itemize}

Если вы хотите \textbf{самостоятельно реализовать Deep Q-Learning}, обязательно ознакомьтесь с реализацией CleanRL: \url{https://github.com/vwxyzjn/cleanrl/blob/master/cleanrl/dqn_atari.py}

\CourseFigure{atari-envs.png}{Примеры сред Atari, в которых можно обучать агентов \protect\href{https://huggingface.co/datasets/huggingface-deep-rl-course/course-images/resolve/main/en/unit4/atari-envs.gif}{[оригинал]}}

\begin{center}\rule{0.5\linewidth}{0.5pt}\end{center}

Поздравляем с завершением этой главы!

Если некоторые понятия до сих пор остаются непонятными\ldots{} это совершенно нормально! \textbf{Так было и у меня, и у всех, кто изучал обучение с подкреплением}.

Прежде чем двигаться дальше, уделите время тому, чтобы \textbf{по-настоящему разобраться в материале и попробовать дополнительные задания}. Важно хорошо освоить эти идеи и сформировать прочную основу.

В следующем модуле \textbf{мы познакомимся с \href{https://optuna.org/}{Optuna}}. Одна из самых важных задач глубокого обучения с подкреплением --- подобрать хорошие гиперпараметры обучения. Библиотека Optuna помогает автоматизировать их поиск.

\paragraph{Этот курс создаётся вместе с вами}\label{this-is-a-course-built-with-you}

Мы хотим постоянно улучшать и обновлять курс с учётом ваших отзывов. Если есть предложения, заполните, пожалуйста, \href{https://forms.gle/3HgA7bEHwAmmLfwh9}{эту форму}.

Мы постоянно улучшаем наши руководства, поэтому \textbf{если вы заметили ошибку в ноутбуке}, \href{https://github.com/huggingface/deep-rl-class/issues}{создайте issue в GitHub-репозитории}.

Увидимся во втором бонусном модуле!

\paragraph{Продолжайте учиться, оставайтесь замечательными!}\label{keep-learning-stay-awesome}

\section{Проверка знаний}\label{quiz}

Лучший способ учиться и \href{https://www.coursera.org/lecture/learning-how-to-learn/illusions-of-competence-BuFzf}{избегать иллюзии компетентности} --- \textbf{проверять себя}. Так вы поймёте, \textbf{какие темы необходимо повторить}.

\subsubsection{Вопрос 1. Мы называли Q-обучение табличным методом. Что такое табличные методы?}\label{q1-we-mentioned-q-learning-is-a-tabular-method.-what-are-tabular-methods}

\subsubsection{Вопрос 2. Почему классическое Q-обучение не подходит для решения игр Atari?}\label{q2-why-cant-we-use-a-classical-q-learning-to-solve-an-atari-game}

\begin{itemize}
\tightlist
\item
\textbf{A.} Среды Atari слишком быстрые для Q-обучения.
\item
\textbf{B.} В средах Atari огромное пространство наблюдений, поэтому создавать и обновлять Q-таблицу было бы неэффективно.
\end{itemize}

\subsubsection{Вопрос 3. Почему при использовании изображений в Deep Q-Learning мы объединяем четыре последовательных кадра?}\label{q3-why-do-we-stack-four-frames-together-when-we-use-frames-as-input-in-deep-q-learning}

\subsubsection{Вопрос 4. Какие две фазы есть в Deep Q-Learning?}\label{q4-what-are-the-two-phases-of-deep-q-learning}

\begin{itemize}
\tightlist
\item
\textbf{A.} Сбор опыта (sampling).
\item
\textbf{B.} Перемешивание (shuffling).
\item
\textbf{C.} Повторное ранжирование (reranking).
\item
\textbf{D.} Обучение (training).
\end{itemize}

\subsubsection{Вопрос 5. Зачем в Deep Q-Learning нужна память воспроизведения опыта?}\label{q5-why-do-we-create-a-replay-memory-in-deep-q-learning}

\subsubsection{Вопрос 6. Как работает Double Deep Q-Learning?}\label{q6-how-do-we-use-double-deep-q-learning}

Поздравляем с завершением теста! Если какие-то вопросы оказались сложными, перечитайте соответствующие разделы, чтобы закрепить знания.

\subsection{Ответы на вопросы и пояснения}\label{quiz-solutions-and-explanations}

\subsubsection{Ответ на вопрос 1}\label{q1-solution}

\emph{Табличные методы} применимы к задачам, в которых пространства состояний и действий достаточно малы, чтобы аппроксимированные функции ценности можно было \textbf{представить массивами и таблицами}. Например, \textbf{Q-обучение --- табличный метод}, поскольку для хранения ценностей пар «состояние--действие» используется таблица.

\subsubsection{Ответ на вопрос 2}\label{q2-solution}

\textbf{Правильный ответ:} B.

\textbf{A.} Среды Atari слишком быстрые для Q-обучения.

\textbf{B.} В средах Atari огромное пространство наблюдений, поэтому создавать и обновлять Q-таблицу было бы неэффективно.

\subsubsection{Ответ на вопрос 3}\label{q3-solution}

Мы объединяем кадры, чтобы \textbf{решить проблему нехватки временной информации}: по одному кадру невозможно определить движение. Например, в Pong агент \textbf{не сможет понять направление движения мяча, если получит только один кадр}.

\CourseFigure{temporal-limitation.jpg}{По одному кадру Pong нельзя определить скорость и направление движения мяча \protect\href{https://huggingface.co/datasets/huggingface-deep-rl-course/course-images/resolve/main/en/unit4/temporal-limitation.jpg}{[оригинал]}}

\CourseFigure{temporal-limitation-2.jpg}{Несколько последовательных кадров Pong позволяют определить направление мяча \protect\href{https://huggingface.co/datasets/huggingface-deep-rl-course/course-images/resolve/main/en/unit4/temporal-limitation-2.jpg}{[оригинал]}}

\subsubsection{Ответ на вопрос 4}\label{q4-solution}

\textbf{Правильные ответы:} A и D.

\textbf{A.} Сбор опыта. Пояснение: агент совершает действия и сохраняет наблюдаемые переходы в памяти воспроизведения.

\textbf{B.} Перемешивание.

\textbf{C.} Повторное ранжирование.

\textbf{D.} Обучение. Пояснение: случайным образом выбираем небольшой пакет переходов и обучаемся на нём, обновляя веса градиентным спуском.

\subsubsection{Ответ на вопрос 5}\label{q5-solution}

\textbf{1. Более эффективное использование опыта во время обучения.}

При обычном онлайн-обучении с подкреплением агент взаимодействует со средой, получает переходы (состояние, действие, вознаграждение и следующее состояние), учится на них, обновляя нейронную сеть, а затем отбрасывает. Это неэффективно. При Experience Replay \textbf{создаётся буфер, в котором хранятся переходы, пригодные для повторного использования во время обучения}.

\textbf{2. Предотвращение забывания и уменьшение корреляции между переходами.}

При последовательной подаче переходов нейронная сеть \textbf{склонна забывать прежний опыт по мере запоминания нового}. Например, если после первого уровня агент попадает на сильно отличающийся второй, он может забыть, как правильно играть на первом уровне.

\subsubsection{Ответ на вопрос 6}\label{q6-solution}

При вычислении целевого Q-значения мы используем две сети, чтобы разделить выбор действия и формирование целевого значения. Для этого мы:

\begin{itemize}
\item
Используем \emph{основную сеть DQN}, чтобы \textbf{выбрать лучшее действие для следующего состояния} (действие с максимальным Q-значением).
\item
Используем \emph{целевую сеть}, чтобы \textbf{рассчитать целевое Q-значение выбранного действия в следующем состоянии}.
\end{itemize}

\section{Заключение}\label{conclusion}

Поздравляем с завершением главы! В ней было много информации. Поздравляем также с завершением практического занятия: вы только что обучили своего первого агента Deep Q-Learning и опубликовали его в Hub.

Прежде чем двигаться дальше, выделите время на то, чтобы хорошо усвоить материал.

Не бойтесь обучать агента в других средах (Pong, Seaquest, QBert, Ms Pac Man). \textbf{Лучший способ научиться --- пробовать всё самостоятельно!}

\CourseFigure{atari-envs.png}{Примеры сред Atari, в которых можно обучать агентов \protect\href{https://huggingface.co/datasets/huggingface-deep-rl-course/course-images/resolve/main/en/unit4/atari-envs.gif}{[оригинал]}}

В следующем модуле \textbf{мы познакомимся с Optuna}. Одна из важнейших задач глубокого обучения с подкреплением --- подбор хорошего набора гиперпараметров. Библиотека Optuna помогает автоматизировать их поиск.

Нам также очень важно \textbf{услышать, что вы думаете о курсе и как его можно улучшить}. Если есть отзывы, пожалуйста, \href{https://forms.gle/BzKXWzLAGZESGNaE9}{заполните эту форму}.

\subsubsection{Продолжайте учиться и оставайтесь замечательными!}\label{keep-learning-stay-awesome-1}

\section{Глоссарий}\label{glossary}

Этот глоссарий создан сообществом. Мы рады новым предложениям и дополнениям!

\begin{itemize}
\item
\textbf{Табличный метод (Tabular Method):} подход для задач с достаточно небольшими пространствами состояний и действий, позволяющими представлять аппроксимированные функции ценности массивами и таблицами. Например, \textbf{Q-обучение} --- табличный метод, так как для хранения ценности различных пар «состояние--действие» используется таблица.
\item
\textbf{Глубокое Q-обучение (Deep Q-Learning):} метод, обучающий нейронную сеть аппроксимировать \textbf{Q-значения} всех возможных действий по текущему состоянию. Он применяется, когда пространство наблюдений слишком велико для табличного Q-обучения.
\item
\textbf{Проблема временной информации (Temporal Limitation):} трудность, возникающая, когда состояние среды представлено отдельными изображениями. Один кадр не передаёт информацию о движении. Чтобы получить её, необходимо \textbf{объединить в стек} несколько последовательных кадров.
\item
\textbf{Фазы Deep Q-Learning:}

  \begin{itemize}
  \tightlist
  \item
\textbf{Сбор опыта (Sampling):} агент совершает действия, а полученные переходы сохраняются в \textbf{буфере воспроизведения опыта}.
  \item
\textbf{Обучение (Training):} случайным образом выбираются пакеты переходов, и нейросеть обновляет веса с помощью градиентного спуска.
  \end{itemize}
\item
\textbf{Способы стабилизации Deep Q-Learning:}

  \begin{itemize}
  \item
\textbf{Воспроизведение опыта (Experience Replay):} создаётся память переходов, которые можно повторно использовать во время обучения. Благодаря этому агент многократно учится на одном и том же опыте. Кроме того, это помогает не забывать прежние переходы по мере поступления новых.
  \item
\textbf{Случайная выборка} из буфера уменьшает корреляцию между последовательными наблюдениями и предотвращает сильные колебания или расходимость ценностей действий.
  \item
\textbf{Фиксированная целевая Q-сеть (Fixed Q-Target):} для вычисления \textbf{целевого Q-значения} с помощью уравнения Беллмана необходимо оценить дисконтированную оптимальную \textbf{Q-ценность} следующего состояния. Проблема в том, что для вычисления \textbf{Q-Target} и текущей \textbf{Q-оценки} используются одни и те же веса сети. Поэтому при каждом изменении \textbf{Q-оценки} смещается и \textbf{цель}. Чтобы избежать этого, для оценки целевого значения временных разностей применяется отдельная сеть с фиксированными параметрами. Её обновляют, копируя веса основной DQN через определённое число \textbf{шагов \(C\)}.
  \item
\textbf{Double DQN:} метод борьбы с \textbf{переоценкой Q-значений}. В нём применяются две сети, позволяющие разделить выбор действия и вычисление целевого \textbf{значения ценности}:

    \begin{itemize}
    \tightlist
    \item
\textbf{Основная сеть DQN} выбирает лучшее действие в следующем состоянии (с максимальным \textbf{Q-значением}).
    \item
\textbf{Целевая сеть} вычисляет целевое \textbf{Q-значение} выбранного действия в следующем состоянии. Такой подход уменьшает переоценку Q-значений, ускоряет обучение и делает его более устойчивым.
    \end{itemize}
  \end{itemize}
\end{itemize}

Если вы хотите улучшить курс, можете \href{https://github.com/huggingface/deep-rl-class/pulls}{создать Pull Request}.

Создание этого глоссария стало возможным благодаря:

\begin{itemize}
\tightlist
\item
\href{https://github.com/dario248}{Дарио Паэсу (Dario Paez)}.
\end{itemize}

\section{Дополнительная литература}\label{additional-readings}

Эти материалы \textbf{необязательны}, но помогут глубже разобраться в теме.

\begin{itemize}
\tightlist
\item
\href{https://youtu.be/Psrhxy88zww}{Серия Foundations of Deep RL, лекция 2 «Deep Q-Learning», Питер Аббил}
\item
\href{https://arxiv.org/abs/1312.5602}{Игра в Atari с глубоким обучением с подкреплением (Playing Atari with Deep Reinforcement Learning)}
\item
\href{https://papers.nips.cc/paper/2010/hash/091d584fced301b442654dd8c23b3fc9-Abstract.html}{Двойное глубокое Q-обучение (Double Deep Q-Learning)}
\item
\href{https://arxiv.org/abs/1511.05952}{Приоритизированное воспроизведение опыта (Prioritized Experience Replay)}
\item
\href{https://arxiv.org/abs/1511.06581}{Дуэльное глубокое Q-обучение (Dueling Deep Q-Learning)}
\end{itemize}

\end{document}
